\section{Motivation}

Environmental, Social, and Governance disclosure has undergone a fundamental paradigm shift, transitioning from a niche voluntary practice to a critical institutional requirement for investors, regulators, and researchers evaluating corporate sustainability performance  ~\cite{hans_bonde_christensen_45b44695}. This evolution reflects a broader commitment to stakeholders beyond traditional shareholders, as firms face mounting pressure to demonstrate tangible "good" alongside financial outcomes  ~\cite{hans_bonde_christensen_45b44695}. However, the current landscape of sustainability reporting remains characterized by extreme structural and content heterogeneity. Reports vary significantly in format and depth due to the absence of global regulatory standardization, often conflating concrete performance metrics with vague, forward-looking narratives  ~\cite[]{ellen_pei_yi_yu_4668244e, hans_bonde_christensen_45b44695}.

In the Indonesian context, these challenges are acutely pronounced. Indonesia is currently classified as having an early-stage reporting framework characterized by lower ESG scores compared to developed markets  ~\cite{monica_singhania_286966ff}. Listed companies often produce bilingual reports with varying degrees of transparency, yet empirical evidence indicates that the readability of these disclosures remains consistently low  ~\cite{desi_adhariani_toit_543335d6}. Linguistic analysis suggests that such reports are often difficult for users to decipher, frequently exhibiting a pattern of "isomorphism" where companies adopt standardized language constructs that obscure specific sustainability issues  ~\cite{desi_adhariani_toit_543335d6}. This lack of clarity facilitates a gap between disclosed intentions and actual corporate conduct, as communication methods remain insufficient to bridge the divide between reporting quality and environmental risk  ~\cite[]{marsdenia_010e13e7, iman_harymawan_2ef898fe}.

Traditional document-level analysis—which assigns a singular score or sentiment to an entire report—is increasingly inadequate for capturing this complexity. ESG disclosures are inherently multifaceted; a single page may contain a specific environmental commitment, a social program update, and a generic governance compliance statement (Mengaldo 2025a). This internal variance allows for "greenwashing," where companies deploy "soft language" or boilerplate narratives to mask a lack of measurable sustainability outcomes  ~\cite{habiba_al_shaer_c2f5394b}. To address these limitations, this research argues that ESG disclosures must be treated as a record-level \textbf{Aspect-Based Sentiment Analysis} problem. By decomposing reports into discrete statements classified by aspect, pillar, sentiment, and provenance, we can transform raw PDF corpora into aligned, inspectable, and evaluable computational artifacts.
\section{Problem Statement}

The central research problem is the absence of a reproducible and validated pipeline capable of transforming heterogeneous sustainability-report PDFs into structured, record-level evidence. Current manual review processes are unscalable, while standard automated tools often fail to capture the semantic nuance and domain-specific logic required for rigorous ESG analysis  ~\cite{p__raghavendra_rau_a32af285}. This study addresses six critical technical challenges:
\begin{enumerate}
\item 

\textbf{Extraction Provenance:} Maintaining a verifiable link between extracted text and original document coordinates is essential to move beyond "black-box" reasoning and provide transaction-level evidence required for auditing  ~\cite{steven_paul_nohr_a6874924}.\item 

\textbf{Prompt Configuration:} Designing reliable Large Language Model prompts is critical to address output volatility; task-specific Chain-of-Thought prompting has been shown to significantly reduce hallucinations in long-context disclosures  ~\cite[]{siqi_sun_c8f93141, association_for_artificial_intelligence_2026_86e12eae}.\item 

\textbf{Granular Labeling:} Record-level labeling is necessary to address the "cherry-picking" and "cheap talk" prevalent in voluntary disclosures, as document-level scoring often masks internal variance  ~\cite[]{julia_bingler_268445e3, habiba_al_shaer_c2f5394b}.\item 

\textbf{Construct Validity:} Identifying the divergence between general sentiment and domain-specific climate labels (e.g., ClimateBERT) is necessary to determine if disclosures reflect substantive risk or merely semantic masking  ~\cite[]{julia_bingler_268445e3, julia_bingler_6b484913}.\item 

\textbf{Quality Auditing:} Automated pipelines are susceptible to OCR failure, schema drift, and missing data, particularly in complex, semi-structured tables  ~\cite[]{siqi_sun_c8f93141, jacob_beck_9548a833}.\item 

\textbf{Evidence Visualization:} Converting extraction outputs into interactive dashboards and semantic graphs is required to facilitate stakeholder engagement and align with emerging regulatory requirements like IFRS S1 and S2  ~\cite{ahmad_zaki_49372c30}.
\end{enumerate}
\section{Research Questions}

The study systematically addresses the gap between raw disclosure and actionable intelligence through the following research framework:
\begin{itemize}
\item 

\textbf{RQ1. ESG ABSA Schema:} How can ESG disclosures be effectively represented using a record-level schema that integrates aspect, ESG pillar, sentiment, and tone within the Indonesian reporting context?\item 

\textbf{RQ2. Tone vs. Climate-Specific Models:} How do LLM-generated tone labels compare with specialized outputs from models like ClimateBERT, and what does this reveal about the validity and reliability of automated ESG assessments?\item 

\textbf{RQ3. Pipeline Diagnostics:} What specific failure modes, such as OCR-related data loss or ontology gaps, characterize the automated extraction of ESG records from heterogeneous PDF sources?\item 

\textbf{RQ4. Stability and Reproducibility:} To what extent do ESG extraction outputs remain stable and consistent across varying prompts, LLM models, and service providers?
\end{itemize}
\section{Objectives and Contributions}

The study systematically addresses the gap between raw disclosure and actionable intelligence through the following research framework:
\subsection{Research Objectives}
\begin{itemize}
\item 

\textbf{O1: To Build a Reproducible OCR-to-ESG Extraction Pipeline:} Develop a modular system that converts raw PDF reports into usable text while preserving page-level provenance for rigorous auditability.\item 

\textbf{O2: To Design a Record-Level ESG ABSA Schema:} Establish a technical framework for classifying extracted statements based on granular ESG pillars and sentiment dimensions tailored to bilingual Indonesian disclosures.\item 

\textbf{O3: To Evaluate Model Divergence:} Conduct a comparative analysis between general LLM tone labels and specialized ClimateBERT labels to assess construct validity and the necessity of domain-specific modeling.\item 

\textbf{O4: To Create a Diagnostic Framework:} Implement automated audits for parsing failures, schema drift, and missing labels to ensure the integrity of the extraction pipeline.\item 

\textbf{O5: To Develop a Semantic Evidence Layer:} Map extracted aspects to a formal ESG ontology and provide interactive dashboards for real-time data exploration and evidence verification  ~\cite{ahmad_zaki_49372c30}.
\end{itemize}
\subsection{Expected Contributions}
\begin{itemize}
\item 

\textbf{Technical Framework:} An executable, modular pipeline for automated ESG record extraction, addressing the current scholarly gap in engaging with quantified sustainability metrics  ~\cite{ahmad_zaki_49372c30}.\item 

\textbf{Annotation Schema:} A bilingual, record-level ESG schema specifically refined for the Indonesian regulatory and linguistic context  ~\cite{desi_adhariani_toit_543335d6}.\item 

\textbf{Evidence Layer:} A linked data structure connecting raw report artifacts to LLM metadata, human-in-the-loop annotations, and validation tables.\item 

\textbf{Methodological Insights:} A comprehensive stability analysis of LLM prompts and a critical assessment of why general sentiment and climate-specific commitment labels are non-interchangeable in high-stakes sustainability research.
\end{itemize}
\section{Thesis Structure}
\begin{itemize}
\item 

\textbf{Chapter 1:} Introduction—Outlines the motivation, problem statement, and technical objectives of the study within the context of global and Indonesian ESG reporting.\item 

\textbf{Chapter 2:} Literature Review—Examines the evolution of NLP in sustainability, the challenge of greenwashing, and the state-of-the-art in LLM-based extraction and ESG measurement discrepancies  ~\cite{p__raghavendra_rau_a32af285}.\item 

\textbf{Chapter 3:} Methodology—Details the OCR-to-ABSA pipeline architecture, including prompt engineering strategies and the validation framework utilizing ClimateBERT and human annotation.\item 

\textbf{Chapter 4:} Implementation and Evaluation—Presents technical implementation details, performance metrics, and comparative diagnostics regarding model stability across different LLM providers.\item 

\textbf{Chapter 5:} Results and Discussion—Evaluates the proposed ESG ABSA schema, analyzes specific failure modes, and discusses the implications of model divergence for future ESG research and policy.\item 

\textbf{Chapter 6:} Conclusion—Summarizes the findings, acknowledges limitations, and suggests future directions for automated, granular sustainability analysis.
\end{itemize}

\section{Author's contributions and the role of Artificial Intelligence}

The author independently formulated the research questions, designed and implemented the framework, and evaluation in the form of live dashboard. The implementation, including dataset preprocessing, software development, and evaluation, was executed by the author under the guidance of the thesis supervisor. 

The writing, including the articulation of arguments and interpretation of findings, reflects the author’s original contributions. Artificial intelligence (AI) tools were employed to enhance the presentation and documentation of this thesis while preserving the integrity of original research
contributions. Specifically, Codex (developed by OpenAI) assisted in refining the clarity
and readability of written content across multiple drafts, as well as provided support for debugging code snippets, resolving LaTeX formatting issues, and structuring tabular data. All AI-generated suggestions underwent rigorous evaluation and extensive revision by the author to ensure alignment with academic standards and research objectives. The fundamental research contributions, including the conceptual framework development, experimental design, data analysis, and scholarly conclusions, represent entirely original work by the author, with AI tools serving exclusively as auxiliary resources for writing enhancement and documentation support.

% \section{Motivation}

% Environmental, Social, and Governance (ESG) disclosure has transitioned from a niche voluntary practice to a critical requirement for investors, regulators, and researchers seeking to evaluate corporate sustainability performance  ~\cite{hans_bonde_christensen_45b44695}. However, the current landscape of sustainability reporting is characterized by extreme heterogeneity. Reports vary significantly in format and extent due to a lack of global standardization, often blending concrete performance data with vague, forward-looking narratives  ~\cite{habiba_al_shaer_c2f5394b}.

% In the Indonesian context, these challenges are further compounded. Listed companies often produce bilingual reports with varying degrees of transparency, and empirical evidence suggests that the readability of Indonesian sustainability reports remains low  ~\cite{desi_adhariani_543335d6}. This lack of clarity often obscures specific sustainability issues, creating a gap between disclosed intentions and actual corporate conduct  ~\cite{marsdenia_010e13e7}.

% Traditional document-level analysis—which assigns a single score or sentiment to an entire report—is increasingly viewed as insufficient. ESG content is inherently multifaceted; a single page might contain an environmental commitment, a social program update, and a generic governance compliance statement  ~\cite{keane_ong_0a5de439}. This complexity allows for "greenwashing," where companies use "soft language" to mask a lack of measurable outcomes  ~\cite{keane_ong_8fec4760}. To address this, this research argues that ESG disclosures must be treated as a record-level \textbf{Aspect-Based Sentiment Analysis} problem. By treating every individual statement as a discrete record classified by aspect, pillar, sentiment, and provenance, we can transform raw PDF corpora into aligned, inspectable, and evaluable computational artifacts.
% \section{Problem Statement}

% The central research problem is the absence of a reproducible and validated pipeline for transforming heterogeneous sustainability-report PDFs into structured, record-level evidence. Current manual review processes are unscalable, while standard automated tools often fail to capture the semantic nuance required for ESG analysis.

% This research addresses six critical technical challenges:
% \begin{enumerate}
% \item 

% \textbf{Extraction Provenance:} Maintaining a clear link between extracted text and its original page/document location in heterogeneous PDFs.\item 

% \textbf{Prompt Configuration:} Designing Large Language Model prompts that can reliably extract multi-faceted ESG statements across different providers.\item 

% \textbf{Granular Labeling:} Assigning precise aspect, pillar, and tone labels at the record level rather than the document level.\item 

% \textbf{Construct Validity:} Identifying the divergence between general sentiment/tone and domain-specific climate labels (e.g., ClimateBERT).\item 

% \textbf{Quality Auditing:} Systematically identifying failures in OCR, schema drift, and missing data fields.\item 

% \textbf{Evidence Visualization:} Converting large-scale extraction results into interactive dashboards and semantic graphs for research verification.
% \end{enumerate}
% \section{Research Questions}

% The study is guided by the following research questions:
% \begin{itemize}
% \item 

% \textbf{RQ1. ESG ABSA Schema:} How can ESG disclosures be effectively represented using a record-level schema that integrates aspect, ESG pillar, sentiment, and tone?\item 

% \textbf{RQ2. Tone vs. Climate-Specific Models:} How do LLM-generated tone labels compare with specialized outputs from models like ClimateBERT, and what does this reveal about the validity of automated ESG assessments?\item 

% \textbf{RQ3. Pipeline Diagnostics:} What specific failure modes, such as OCR-related data loss or ontology gaps, characterize the automated extraction of ESG records?\item 

% \textbf{RQ4. Stability and Reproducibility:} To what extent do ESG extraction outputs remain stable across varying prompts, LLM models, and service providers?
% \end{itemize}
% \section{Objectives and Contributions}
% \subsection{Research Objectives:}
% \begin{itemize}
% \item 

% \textbf{O1: To Build a Reproducible OCR-to-ESG Extraction Pipeline:} Develop a system that converts raw PDF reports into usable text while preserving page-level provenance for auditability.\item 

% \textbf{O2: To Design a Record-Level ESG ABSA Schema:} Establish a technical framework for classifying extracted statements based on granular ESG pillars and sentiment dimensions.\item 

% \textbf{O3: To Evaluate Model Divergence:} Conduct a comparative analysis between general LLM tone labels and specialized ClimateBERT climate labels to assess construct validity.\item 

% \textbf{O4: To Create a Diagnostic Framework:} Implement automated audits for parsing failures, schema drift, and missing labels within the extraction pipeline.\item 

% \textbf{O5: To Develop a Semantic Evidence Layer:} Map extracted aspects to a formal ESG ontology and provide interactive dashboards for real-time data exploration and evidence verification.
% \end{itemize}
% \subsection{Expected Contributions:}
% \begin{itemize}
% \item 

% \textbf{Technical Framework:} An executable, modular pipeline for automated ESG record extraction from sustainability reports.\item 

% \textbf{Annotation Schema:} A bilingual, record-level ESG schema tailored for the Indonesian reporting context.\item 

% \textbf{Evidence Layer:} A linked data structure connecting raw reports to LLM metadata, human annotations, and validation tables.\item 

% \textbf{Methodological Insights:} A comprehensive stability analysis of LLM prompts and a critical discussion on why general sentiment and climate-specific commitment labels are non-interchangeable.
% \end{itemize}
% \section{Thesis Structure}
% \begin{itemize}
% \item 

% \textbf{Chapter 1:} Outlines the motivation, problem statement, and technical objectives of the study.\item 

% \textbf{Chapter 2:} Examines the evolution of NLP in sustainability, the challenge of greenwashing, and the current state-of-the-art in LLM-based extraction.\item 

% \textbf{Chapter 3:} Details the OCR-to-ABSA pipeline, including prompt engineering strategies and the validation framework using ClimateBERT and human annotation.\item 

% \textbf{Chapter 4:} Presents implementation details, evaluation metrics, and ablation studies regarding model stability.\item 

% \textbf{Chapter 5:} Evaluates the ESG ABSA schema, analyzes failure modes, and discusses the implications of model divergence for future ESG research.\item 

% \textbf{Chapter 6:} Summarizes the findings and suggests future directions for automated sustainability analysis.
% \end{itemize}
